\section{Implementación a nivel de la API de OpenAI}

\subsection{Structured Output: el campo response\_format}

Se pasa el JSON Schema del modelo de Pydantic a través del campo \texttt{response\_format}:

\begin{itemize}
    \item Se define un modelo de Pydantic (por ejemplo, con los campos title y year)
    \item Se serializa a JSON Schema mediante \texttt{model.json\_schema()}
    \item Se coloca en el campo \texttt{response\_format} de la solicitud HTTP
    \item La salida del modelo se ajusta automáticamente a ese Schema y puede deserializarse directamente en un objeto
\end{itemize}

\subsection{Function Calling: el campo tools}

Se inyecta la información de las herramientas mediante el campo \texttt{tools}:

\begin{itemize}
    \item Se registra la función de Python con el decorador \texttt{@tool} de LangChain
    \item Se convierte al formato que exige la API de OpenAI mediante \texttt{convert\_to\_openai\_tool()}
    \item Se coloca en el campo \texttt{tools} de la solicitud HTTP
    \item La respuesta del modelo incluye el contenido instanciado de \texttt{function\_call}
\end{itemize}

\begin{knowledgebox}{Análisis por captura de tráfico HTTP}
Mediante la captura de tráfico HTTP se puede ver con claridad que:
\begin{itemize}
    \item \textbf{Structured Output} pasa por el campo \texttt{response\_format}
    \item \textbf{Function Calling / Bound Tools} pasa por el campo \texttt{tools}
\end{itemize}
Ambos son campos distintos a nivel de la solicitud HTTP, pero su propósito es el mismo: fijar el formato de entrada y salida.
\end{knowledgebox}

\subsection{Resumen de la sección}

A nivel de la solicitud HTTP de la API de OpenAI, la salida estructurada se implementa mediante response\_format y la llamada a herramientas mediante tools; entender el mecanismo subyacente ayuda a depurar y optimizar.

